\normaltrue \difficilefalse \tdifficilefalse
\correctiontrue

%\UPSTIidClasse{11} % 11 sup, 12 spé
%\newcommand{\UPSTIidClasse}{11}

\exer{Diagramme de Bode$\star$ \label{C2:02:510_01}}
\setcounter{question}{0}\marginnote{\xpComp{SLCI}{11}}%\UPSTIcompetence[2]{C2-02}
\index{Compétence C2-02}\index{Compétence SLCI-11}
\index{Diagramme de Bode}
\ifcorrection
\else
\marginnote{\textbf{Pas de corrigé pour cet exercice.}}
\fi

\question{Tracer le diagramme de Bode de la fonction de transfert suivante : $F_1(p)=\dfrac{15}{1+10p}$.}

\ifprof
\textbf{Tracé asymptotique}

\begin{marginfigure}

\includegraphics[width=\linewidth]{tab_01}
\end{marginfigure}

\textbf{Positionnement du diagramme de gain}
Lorsque $\omega$ tend vers 0, le gain tend vers
$20\log 15=\SI{23,5}{dB}$.

% Le tracé historique utilisait bodegraph/gnuplot et nécessitait
% la génération externe d'un fichier .table. Le PDF déjà présent
% dans la banque est utilisé directement pour garantir une compilation
% autonome de ALL_EXOS sous MiKTeX.
\begin{figure}[!h]
\centering
\includegraphics[width=\linewidth]{bode_01}
\end{figure}
\else
\fi

\question{Le système est sollicité par une entrée sinusoïdale de période \SI{60}{s} et d'amplitude 10. Quel est le signal de sortie ?}
\ifprof
Pour une période de \SI{60}{s}, la pulsation est de $\dfrac{2\pi}{T}$,
soit $\omega=\SI{0,1}{rad.s^{-1}}$.
Pour cette pulsation, le gain est de \SI{20}{dB} et le déphasage de
$-\dfrac{\pi}{4}$.

On a donc $20\log(S/E)=20$, soit $S=10E$. Le signal d'entrée est
$e(t)=10\sin(0,1t)$ et le signal de sortie est
$s(t)=100\sin\left(0,1t-\dfrac{\pi}{4}\right)$.
\else
\fi

\ifprof
\else
\marginnote{Corrigé voir \ref{C2:02:510_01}.}
\fi

% BEGIN AUTO CORRIGE INCLUDE
\InclureCorrige{corrige.tex}
% END AUTO CORRIGE INCLUDE
